% Problem record op_cdd1f718ccfbccf6. This TeX file is derived from
% database/problems_json/op_cdd1f718ccfbccf6.json by scripts/sync-tex.mjs; edit the JSON record
% and rerun the script rather than editing this file.
\section{Semi-Clifford structure of two-qudit hierarchy gates above the third level}
\label{sec:op_cdd1f718ccfbccf6}

\paragraph{Problem.}
Is every two-qudit gate in every level $k\geq4$ of the Clifford hierarchy
semi-Clifford when the local dimension is an odd prime?  Let $p$ be an odd
prime, let $\omega:=e^{2\pi i/p}$, and define on $\mathbb{C}^p$ the operators
$X|j\rangle=|j+1\bmod p\rangle$ and $Z|j\rangle=\omega^j|j\rangle$ for
$j\in\{0,\ldots,p-1\}$.  For $n$ qudits, let $\mathcal{P}_{n,p}$ be the group
generated by arbitrary global phases and the single-site operators $X_i,Z_i$
($1\leq i\leq n$), each acting as the identity on the other sites, and define
\begin{equation}
  \mathcal{C}_1(n,p):=\mathcal{P}_{n,p},
  \qquad
  \mathcal{C}_{\ell+1}(n,p):=
  \bigl\{U\in U(p^n):UPU^\dagger\in\mathcal{C}_\ell(n,p)
  \ \text{for every}\ P\in\mathcal{P}_{n,p}\bigr\},
  \qquad \ell\geq1.
  \label{eq:cdd1-hierarchy}
\end{equation}
A unitary is \emph{semi-Clifford} when it equals $C_LDC_R$ with
$C_L,C_R\in\mathcal{C}_2(n,p)$ and $D$ diagonal in the computational basis;
write $\mathrm{SC}(n,p)$ for this set.  The question asks whether, with the
levels of Eq.~\eqref{eq:cdd1-hierarchy},
\begin{equation}
  \mathcal{C}_k(2,p)\subseteq\mathrm{SC}(2,p)
  \qquad\text{for every odd prime }p\text{ and every integer }k\geq4.
  \label{eq:cdd1-inclusion}
\end{equation}
A complete answer either proves Eq.~\eqref{eq:cdd1-inclusion} or exhibits an
odd prime $p$, a level $k\geq4$, and a gate in $\mathcal{C}_k(2,p)$ that is not
semi-Clifford.

\subsection*{Status}

Unsolved

\subsection*{Source}

De Silva conjectured that every $k$th-level gate of one or two qudits of
any prime dimension is semi-Clifford (Conjecture~3 in Section~V.C of the arXiv
version), and in the concluding Section~VI proposed the two-qudit, all-level
statement again by analogy with the two-qubit theorem of Zeng, Chen, and Chuang
\sourcecite{ref:cdd1-ds21}{dS21}.  Chen and de Silva, having proved the
third-level case, name its generalization to higher levels as a natural
follow-up question in their concluding Section~4
\sourcecite{ref:cdd1-cds24}{CdS24}.  Eq.~\eqref{eq:cdd1-inclusion} is the part
of the two-qudit conjecture not settled by the results below: qubits and levels
$k\leq3$ are excluded.

\subsection*{Progress}

\begin{itemize}
  \item For qubits, Zeng, Chen, and Chuang proved by induction on the level,
  using exhaustive computations, that every two-qubit hierarchy gate is
  semi-Clifford (Theorem~1), and showed that the property fails for three qubits
  at every level $k\geq4$ (Theorem~3):
  \begin{equation}
    \mathcal{C}_k(2,2)\subseteq\mathrm{SC}(2,2)\quad(k\geq1),
    \qquad
    \mathcal{C}_k(3,2)\not\subseteq\mathrm{SC}(3,2)\quad(k\geq4),
    \label{eq:cdd1-qubits}
  \end{equation}
  where $\mathcal{C}_k(n,2)$ and $\mathrm{SC}(n,2)$ denote the sets defined as in
  Eq.~\eqref{eq:cdd1-hierarchy} with $p=2$.  The two-qubit case motivates
  Eq.~\eqref{eq:cdd1-inclusion}, but its computational proof does not extend to
  odd prime dimensions \sourcecite{ref:cdd1-zcc08}{ZCC08}.

  \item De Silva proved that every third-level gate of one qudit of any prime
  dimension (Theorem~8) and every third-level gate of two qutrits (Theorem~9, by
  exhaustive computation) is semi-Clifford, and reported numerical evidence that
  the property extends to further cases.  These results do not reach levels
  $k\geq4$ for two qudits \sourcecite{ref:cdd1-ds21}{dS21}.

  \item Chen and de Silva proved that, for every odd prime $p$, every two-qudit
  third-level gate is semi-Clifford (Theorem~1):
  \begin{equation}
    \mathcal{C}_3(2,p)\subseteq\mathrm{SC}(2,p)
    \qquad(p\ \text{an odd prime}).
    \label{eq:cdd1-third}
  \end{equation}
  Their argument describes a simplified two-qudit third-level gate by the
  Clifford gates to which it conjugates the basic Pauli operators, turns the
  third-level and semi-Clifford conditions into polynomial systems over
  $\mathbb{Z}_p$, and compares the rational points of the resulting algebraic
  sets uniformly in $p$.  At levels $k\geq4$ these conjugates are no longer
  Clifford gates, and Eq.~\eqref{eq:cdd1-third} does not establish
  Eq.~\eqref{eq:cdd1-inclusion} \sourcecite{ref:cdd1-cds24}{CdS24}.

  \item De Silva and Lautsch proved that, for every odd prime $p$, every one-qudit
  hierarchy gate is semi-Clifford (Theorem~6), and derived a unique normal form
  $MDC$ for non-Clifford one-qudit hierarchy gates, with $M$ from a fixed finite
  set of Clifford gates, $D$ diagonal, and $C$ Clifford (Theorem~7):
  \begin{equation}
    \mathcal{C}_k(1,p)\subseteq\mathrm{SC}(1,p)
    \qquad(p\ \text{an odd prime},\ k\geq1).
    \label{eq:cdd1-one}
  \end{equation}
  In Section~5 (Conjecture~1 of the arXiv version) they conjecture that every
  two-qudit hierarchy gate is either Clifford or uniquely of the form
  $M_1M_2DC$ with Clifford gates $M_1,M_2$ from explicit finite families, $D$
  diagonal, and $C$ Clifford.  Such a gate is semi-Clifford, so that unproved
  normal form would imply Eq.~\eqref{eq:cdd1-inclusion}; the one-qudit theorem in
  Eq.~\eqref{eq:cdd1-one} does not \sourcecite{ref:cdd1-dsl25}{dSL25}.
\end{itemize}

\subsection*{References}

\begin{enumerate}[label={},leftmargin=4.8em,labelsep=0.5em]
  \footnotesize
  \item[\textup{[ZCC08]}]\label{ref:cdd1-zcc08}
  B. Zeng, X. Chen, and I. L. Chuang, ``Semi-Clifford operations, structure of
  $\mathcal{C}_k$ hierarchy, and gate complexity for fault-tolerant quantum
  computation,'' \emph{Physical Review A} \textbf{77}, 042313 (2008).
  \href{https://doi.org/10.1103/PhysRevA.77.042313}{doi:10.1103/PhysRevA.77.042313};
  \href{https://arxiv.org/abs/0712.2084}{arXiv:0712.2084}.

  \item[\textup{[dS21]}]\label{ref:cdd1-ds21}
  N. de Silva, ``Efficient quantum gate teleportation in higher dimensions,''
  \emph{Proceedings of the Royal Society A} \textbf{477}, 20200865 (2021).
  \href{https://doi.org/10.1098/rspa.2020.0865}{doi:10.1098/rspa.2020.0865};
  \href{https://arxiv.org/abs/2011.00127}{arXiv:2011.00127}.

  \item[\textup{[CdS24]}]\label{ref:cdd1-cds24}
  I. Chen and N. de Silva, ``Characterising semi-Clifford gates using
  algebraic sets,'' \emph{Communications in Mathematical Physics} \textbf{405},
  201 (2024).
  \href{https://doi.org/10.1007/s00220-024-05050-2}{doi:10.1007/s00220-024-05050-2};
  \href{https://arxiv.org/abs/2309.15184}{arXiv:2309.15184}.

  \item[\textup{[dSL25]}]\label{ref:cdd1-dsl25}
  N. de Silva and O. Lautsch, ``The Clifford hierarchy for one qubit or
  qudit,'' \emph{Proceedings of the Royal Society A} \textbf{481}, 20250035
  (2025).
  \href{https://doi.org/10.1098/rspa.2025.0035}{doi:10.1098/rspa.2025.0035};
  \href{https://arxiv.org/abs/2501.07939}{arXiv:2501.07939}.

  \item[\textup{[dSL26]}]\label{ref:cdd1-dsl26}
  N. de Silva and O. Lautsch, ``The generalised semi-Clifford conjecture is
  false,'' arXiv preprint, preliminary draft (2026).
  \href{https://doi.org/10.48550/arXiv.2609.11903}{doi:10.48550/arXiv.2609.11903};
  \href{https://arxiv.org/abs/2609.11903}{arXiv:2609.11903}.
\end{enumerate}

\subsection*{Comment}

What remains open is Eq.~\eqref{eq:cdd1-inclusion} at every level $k\geq4$
for every odd prime $p$.  The established results cover two qubits at every
level, Eq.~\eqref{eq:cdd1-qubits}; two odd-prime-dimensional qudits at level
three, Eq.~\eqref{eq:cdd1-third}; and one qudit at every level,
Eq.~\eqref{eq:cdd1-one}.  None covers two odd-prime-dimensional qudits at levels
four and above.  A counterexample must be a gate in $\mathcal{C}_k(2,p)$ for
the two-qudit Pauli group of Eq.~\eqref{eq:cdd1-hierarchy}; reproducing a
multiqubit counterexample inside an encoded subspace does not qualify.  For
qubits, the weaker generalized semi-Clifford property, which allows an
additional computational-basis permutation between the Clifford factors, fails
for a five-qubit gate in the fifth level, and the same authors announce
forthcoming qudit counterexamples without specifying the number of qudits or
the level \sourcecite{ref:cdd1-dsl26}{dSL26}; a two-qudit gate at some level
$k\geq4$ that is not generalized semi-Clifford would in particular answer the
question negatively.  Literature checked through 15 September 2026; no proof
of Eq.~\eqref{eq:cdd1-inclusion} and no two-qudit counterexample was found.
The zoo's record
``\href{https://qiqc-op.com/problem/op_5a17b19b8adeb191/}{Smallest qubit number for a non-semi-Clifford third-level gate}''
concerns the least number of qubits at which the third level stops being
semi-Clifford.

\subsection*{Field}

Quantum algorithm

\subsection*{Topic}

Clifford hierarchy

\subsection*{ID}

\texttt{op\_cdd1f718ccfbccf6}
